% TOP_PROBLEM: {"id": 2578, "title": "Kourovka Notebook Problem 21.69", "category_id": 17, "category": {"id": 17, "name": "group_theory", "display_name": "Group Theory", "description": "Problems about groups, group actions, representations, and related algebraic structures.", "slug": "group-theory", "order_index": 17, "created_at": "2026-07-31T15:26:25.671Z"}, "status": "open", "problem_number": "KOU-21.69", "set_id": 10}
% FIRST_POSTED: 2026-10-01
% ENTRY_KIND: statement_only
% UnsolvedMath ID 2578; source catalogue code KOU-21.69
% !TeX program = lualatex
% Definition and sources only; this file is not an LLM solution attempt.
\documentclass[11pt,a4paper]{article}
\usepackage[margin=25mm]{geometry}
\usepackage{fontspec}
\setmainfont{lmroman10-regular.otf}[BoldFont=lmroman10-bold.otf,
 ItalicFont=lmroman10-italic.otf,BoldItalicFont=lmroman10-bolditalic.otf]
\usepackage{amsmath,amssymb,mathtools}
\usepackage{unicode-math}
\setmathfont{latinmodern-math.otf}
\AtBeginDocument{\let\setminus\smallsetminus}
\usepackage{xurl,hyperref}
\hypersetup{colorlinks=true,urlcolor=blue}
\setcounter{secnumdepth}{0}
\setlength{\parindent}{0pt}
\setlength{\parskip}{5pt}
\setlength{\emergencystretch}{3em}
\begin{document}
\section*{Kourovka Notebook Problem 21.69}\label{2578}
{\small UnsolvedMath ID 2578 · KOU-21.69 · Group Theory · Kourovka Notebook - New Problems, Issue 21}
\subsection{Definitions and mathematical statement}
A finite one-relator presentation has the form
\[
 G=\langle x_1,\ldots,x_n\mid R\rangle
   =F_n/\langle\!\langle R\rangle\!\rangle,
\]
where $n$ is finite, $F_n$ is the free group on the listed generators,
$R$ is a finite word in them and their inverses, and the double angle
brackets denote the normal subgroup generated by $R$.
The input consists of this finite list and word. No restriction is
imposed on the number of occurrences of a generator in the relator,
and proper-power relators are allowed.

Give the Cayley graph of $G$ for the indicated symmetric generating
set edges of length one and its induced path metric. The group is
word-hyperbolic if some real $\delta\ge0$ has the following property:
every side of every geodesic triangle lies within distance $\delta$
of the union of the other two sides. This property is independent
of the chosen finite generating set.

Is there an algorithm that, for every finite one-relator presentation,
halts and decides whether its presented group is word-hyperbolic?
The algorithm must handle both positive and negative instances; a
procedure that eventually verifies hyperbolicity when it holds but
may run forever otherwise does not suffice.

The question asks for one uniform algorithm over all such presentations,
not a decision procedure tailored to a single relator. No value of
$\delta$ is supplied as part of the input, and no running-time bound
is required. Finite groups and other elementary hyperbolic groups are
included in the positive class.
\subsection{Short English statement}
Can one algorithm decide word-hyperbolicity from every finite one-relator presentation?
\subsection{Sources}
\begin{itemize}
\item[S1] ULAM.AI and the UnsolvedMath Contributors, supplied problems.json, record 2578 (KOU-21.69), Kourovka Notebook - New Problems, Issue 21. Catalogue identity and source transcription; CC BY 4.0.
\url{https://huggingface.co/datasets/ulamai/UnsolvedMath}.
\item[S2] The Kourovka Notebook, 21st edition, new problems, Problem 21.69. Expanded formulation of the supplied catalogue statement.
\url{https://arxiv.org/pdf/1401.0300v47}.
\end{itemize}
\end{document}
